\section{团队结构与运行时干预（RQ2）}
\label{sec:structure}

在单个上下文装得下的任务上，一个能力强的负责者不输于我们搭建的任何团队，成本只是团队的一小部分。增加的智能体在
它们所针对的题上胜出，在不针对的题上持平或落后，成本通常是两到五倍。

\para{小任务不委派，大任务才委派}在第一批任务集上（9 月 28 日），默认团队从不委派，单个智能体同样好，而且略便宜
（10.9 分钟对 12.1 分钟，1.24M token 对 1.41M）。在更大的多模块任务上（9 月 29 日），领队在其中两道题的全部六次运行里
都做了委派，在调研类任务上一次也没有；36 次运行全部解出，所以结构只能影响时间和成本。

\para{为难题搭建的结构不能迁移}从 10 月 8 日起，我们针对上线团队做不出的难题，每天搭一种结构（表~\ref{tab:structures}）。
对这些失败的诊断发现，每一次都是对需求中一句含糊或没说清的话的误读，而且多数运行都以同样方式误读，因此投票和
增加审查者几乎无从下手。在每个节点上加阅读规则，到了样本外反而更差，因为这些规则让模型觉得可以跳过文字没有
明说的部分。

\begin{table}[t]
\centering
\footnotesize
\caption{V1 中为难题搭建的团队结构，每个都与同批基线比较（实验组对基线，通过的运行数）。设计题是该结构所针对的题；
留出题没有用于搭建它。}
\label{tab:structures}
\begin{tabularx}{\textwidth}{@{}>{\raggedright\arraybackslash}p{34mm} L L >{\raggedright\arraybackslash}p{21mm} >{\raggedright\arraybackslash}p{27mm}@{}}
\toprule
结构 & 设计题 & 留出题 & 每次运行成本 & 结论 \\
\midrule
每个节点加阅读规则 & 9/15 对 4/15 & 2/10 对 4/10 & 约 1.00 美元对 0.76 美元 & 放弃：规则让工作缩水 \\
三个构建者加一个证据评审（集成） & 15/24 对 4/15；某道题第一次被解出 & 3/10 对 8/20 & 2.85 美元对 0.53 美元 & 放弃 \\
三个阅读者加一个调和者 & 4/10 对 4/21 & 1/9 对 11/28 & 4.01 美元对 1.31 美元 & 放弃 \\
与构建同时启动的验收检查者 & 3/5 对 2/15（细节密集的题） & 8/10 对 7/10（新题集） & 2.03 美元对 1.58 美元 &
  作为关口采用，只用于规格类需求 \\
画类型化连线图的架构师 & 0/6；五次完全打平 & 未运行 & 未记录 & 放弃 \\
\bottomrule
\end{tabularx}
\end{table}

\para{V1 上线了什么}在留出题上与基线持平的单个检查者确实上线了，一次加一个关口：验收（团队 6.0）、针对缺陷报告的
症状检查（6.1）、代码仓库惯例（6.2）、检查者交回时的定夺步骤（6.3；在其目标题上 17/48 对 10/48）以及组合测试者
（6.4）。更重的变体全都没通过留出检验，包括双检查者版本（3/9 对 5/9）和定夺者节点（合计 20/40 对 21/40）。10 月 18 日对
120 次运行中 61 次失败的普查发现：33 次是误读需求，16 次是自测比隐藏测试宽松，12 次是领域知识缺口。

\para{单领队在 Terminal-Bench 上同样胜出}在 V8 中，团队 7.0（单个领队直接干活，只在被要求时委派）以 max 强度运行，
在未读的 30 道题上比之前中等强度的领队高 $+3.17$（标准误 1.39），10 道题上升、3 道下降，没有丢掉任何一道原本总能
通过的题，成本是两倍。这是两个配置之间的比较：领队的人设和推理强度一起变了，而这一对组合正是我们上线的
（单独比较推理强度见第~\ref{sec:judgment}~节）。按 Codex 自己的协作协议做主动委派（领队不等要求就派生助手智能体），
只在 24 次目标试验中的 2 次上成功，丢了两道守卫题，成本是二到四倍；它的被动形式，即只在被要求时委派，就是团队 7.0
采用的方式。一个检查型团队与单领队持平，成本是两倍。

\para{同模型的第二意见在样本外无效}在动手之前先研究题目的阅读者，在未读题上合计为 $-0.8$（$\pm 0.7$）。在两份已完成
的构建之间做选择的选择者，无论用 GPT 还是 Claude，正确率都在 55\% 左右：失败的那份有明显缺陷时几乎总能选对
（Claude Opus 23 次中 22 次），在需要判断的点上则很少选对（28 次中 5 次）。一个同模型陪审团在 36 对“通过/失败”构建上
判对 12 次、判错 11 次、弃权 13 次。由领队启动的独立评分者在目标试验上通过 3/10，基线为 5/10；它给出的 10 个
“全部通过”结论中有 6 个落在失败的试验上。一个带关口的推理检查者以 12/20 打平。

\para{框架只在回应“通过的运行从不出现的状态”时有用}无界面运行中没有人回答交回。当一次运行很早就交回失败时，框架
回答“你自己决定，继续做”（GO）；有一道题从 14 次中 0 次通过变成 6 次中 6 次，GO 因此上线。在每次首次交付时触发的
检查没有上线。GPT-6.1 上的 HD 在未读题上得 $+1.54$（$\pm 1.31$），成本 2.09 倍，低于预先登记的 $+2.0$ 门槛。O2 是同一
思路针对 GPT-6.2 的措辞，在多次运行结果不一致的可读题上通过 22/28 对 17/28，它的第二轮确实找出了真缺陷；但在未读题上
少了 2.0（15.5 对 17.5，$p = 0.28$），并在 60 次试验中的 50 次里返工了已交付的工作。全局触发的手段在一个失败类型里
得到的收益，要用原本已经做对的题来偿还。

\para{RQ2 的回答}在单上下文任务上，同一模型组成的团队增加了成本，没有提高准确率。有用的东西都很窄：一个只用于
规格类需求的关口，以及在运行时回应一种通过的运行从不会进入的状态（很早就交回失败）。
